% GENERATED FILE. Edit canonical lessons, then run npm run generate:books.
% canonical-source-hash: fnv1a64:5fe81b92b89f787d
% canonical-lessons: ZH-C202-saoma, ZH-C202-shouju, ZH-C202-tuihuo, ZH-C202-guke, ZH-C202-shouyintai

\chapter{At the Till}
\label{ch:zh-at-the-till}
\hlchaptermodality{202}

\begin{chapteropening}
\textbf{By the end of this chapter:} \emph{I can pay by scanning a code, ask for a receipt, take something back and find the checkout.}

Complete the last lesson of chapter 202: I can pay by scanning a code, ask for a receipt, take something back and find the checkout.
\end{chapteropening}

\section[sǎomǎ]{sǎomǎ — to scan a code (to pay)}

\label{lesson:ZH-C202-saoma}



\begin{quote}
Before the new one: say the Mandarin for to send electronically, then the Mandarin for to get something sent.
\end{quote}

\subsection*{You'll want to know: sǎomǎ}
\emph{sǎomǎ} — \textquotedblleft{}to scan a code (to pay)\textquotedblright{}.

\begin{quote}
\emph{kěyǐ sǎomǎ ma} — Can I pay by scanning the code?
\end{quote}

Its characters come later; for now it is a word for the ear.

Two third tones meet in this word: sǎo is said with a rising tone before mǎ, and the spelling keeps both marks.

\begin{cousinweb}
\emph{sǎomǎ} is built from \zh{扫} 'to sweep' and \emph{mǎ} 'a code', both inherited from Middle Chinese, \emph{sawX} and \emph{maeX} (Guangyun, in Baxter's transcription). The pair is recent, and how sweeping came to mean scanning is not settled.
\end{cousinweb}

\begin{grammarlens}[title={What you've built}]
One more word for asking for what I want.
\end{grammarlens}

\subsection*{Your turn}
\begin{itemize}
\raggedright
  \item \emph{Say it:} \emph{sǎomǎ}
  \item \emph{Say it:} \emph{sǎomǎ}, once more
  \item \emph{Say it:} \emph{kěyǐ sǎomǎ ma} — Can I pay by scanning the code?
  \item \emph{Recall:} say the Mandarin for to send electronically, then the Mandarin for to get something sent, then say \emph{sǎomǎ} again
\end{itemize}

\begin{tcolorbox}[breakable,colback=teal!4,colframe=teal!35!black,title={Before you move on}]
What does sǎomǎ mean? (\textquotedblleft{}To scan a code (to pay)\textquotedblright{}.) Say it once more.
\end{tcolorbox}
\section[shōujù]{shōujù — receipt}

\label{lesson:ZH-C202-shouju}



\begin{quote}
Before the new one: say the Mandarin for to get something sent, then the Mandarin for to scan a code.
\end{quote}

\subsection*{You'll want to know: shōujù}
\emph{shōujù} — \textquotedblleft{}receipt\textquotedblright{}.

\begin{quote}
\emph{wǒ yào yì zhāng shōujù} — I need a receipt.
\end{quote}

Its characters come later; for now it is a word for the ear.

\begin{cousinweb}
\emph{shōujù} is built from \emph{shōu} 'to receive' and \emph{jù} 'evidence', both inherited from Middle Chinese, \emph{syuw} and \emph{kjoH} (Guangyun, in Baxter's transcription), which glosses the second 'to rely on'.
\end{cousinweb}

\begin{grammarlens}[title={What you've built}]
One more word for asking for what I want.
\end{grammarlens}

\subsection*{Your turn}
\begin{itemize}
\raggedright
  \item \emph{Say it:} \emph{shōujù}
  \item \emph{Say it:} \emph{shōujù}, once more
  \item \emph{Say it:} \emph{wǒ yào yì zhāng shōujù} — I need a receipt.
  \item \emph{Recall:} say the Mandarin for to get something sent, then the Mandarin for to scan a code, then say \emph{shōujù} again
\end{itemize}

\begin{tcolorbox}[breakable,colback=teal!4,colframe=teal!35!black,title={Before you move on}]
What does shōujù mean? (\textquotedblleft{}Receipt\textquotedblright{}.) Say it once more.
\end{tcolorbox}
\section[tuìhuò]{tuìhuò — to take goods back; to return goods to a shop}

\label{lesson:ZH-C202-tuihuo}



\begin{quote}
Before the new one: say the Mandarin for to scan a code, then the Mandarin for receipt.
\end{quote}

\subsection*{You'll want to know: tuìhuò}
\emph{tuìhuò} — \textquotedblleft{}to take goods back; to return goods to a shop\textquotedblright{}.

\begin{quote}
\emph{zhège kěyǐ tuìhuò ma} — Can I bring this back?
\end{quote}

Its characters come later; for now it is a word for the ear.

\begin{cousinweb}
\emph{tuìhuò} is built from \emph{tuì} 'to give back' and \emph{huò} 'goods', both inherited from Middle Chinese, \emph{thwojH} and \emph{xwaH} (Guangyun, in Baxter's transcription), glossed 'to draw back' and 'wealth'.
\end{cousinweb}

\begin{grammarlens}[title={What you've built}]
One more word for asking for what I want.
\end{grammarlens}

\subsection*{Your turn}
\begin{itemize}
\raggedright
  \item \emph{Say it:} \emph{tuìhuò}
  \item \emph{Say it:} \emph{tuìhuò}, once more
  \item \emph{Say it:} \emph{zhège kěyǐ tuìhuò ma} — Can I bring this back?
  \item \emph{Recall:} say the Mandarin for to scan a code, then the Mandarin for receipt, then say \emph{tuìhuò} again
\end{itemize}

\begin{tcolorbox}[breakable,colback=teal!4,colframe=teal!35!black,title={Before you move on}]
What does tuìhuò mean? (\textquotedblleft{}To take goods back; to return goods to a shop\textquotedblright{}.) Say it once more.
\end{tcolorbox}
\section[gùkè]{gùkè — customer}

\label{lesson:ZH-C202-guke}



\begin{quote}
Before the new one: say the Mandarin for receipt, then the Mandarin for to take goods back.
\end{quote}

\subsection*{You'll want to know: gùkè}
\emph{gùkè} — \textquotedblleft{}customer\textquotedblright{}.

\begin{quote}
\emph{zhèr de gùkè dōu hěn niánqīng} — The customers here are all young.
\end{quote}

Its characters come later; for now it is a word for the ear.

\begin{cousinweb}
\emph{gùkè} is built from \emph{gù} 'to look round' and \emph{kè} 'a guest', both inherited from Middle Chinese, \emph{kuH} and \emph{khaek} (Guangyun, in Baxter's transcription), which glosses the first 'to look back'.
\end{cousinweb}

\begin{grammarlens}[title={What you've built}]
One more word for asking for what I want.
\end{grammarlens}

\subsection*{Your turn}
\begin{itemize}
\raggedright
  \item \emph{Say it:} \emph{gùkè}
  \item \emph{Say it:} \emph{gùkè}, once more
  \item \emph{Say it:} \emph{zhèr de gùkè dōu hěn niánqīng} — The customers here are all young.
  \item \emph{Recall:} say the Mandarin for receipt, then the Mandarin for to take goods back, then say \emph{gùkè} again
\end{itemize}

\begin{tcolorbox}[breakable,colback=teal!4,colframe=teal!35!black,title={Before you move on}]
What does gùkè mean? (\textquotedblleft{}Customer\textquotedblright{}.) Say it once more.
\end{tcolorbox}
\section[shōuyíntái]{shōuyíntái — checkout; till}

\label{lesson:ZH-C202-shouyintai}



\begin{quote}
Before the new one: say the Mandarin for to take goods back, then the Mandarin for customer.
\end{quote}

\subsection*{You'll want to know: shōuyíntái}
\emph{shōuyíntái} — \textquotedblleft{}checkout; till\textquotedblright{}.

\begin{quote}
\emph{shōuyíntái zài ménkǒu} — The checkout is by the door.
\end{quote}

Its characters come later; for now it is a word for the ear.

\begin{cousinweb}
\emph{shōuyíntái} is built from \emph{shōu} 'to take in', \emph{yín} 'silver, money' and \emph{tái} 'a counter', all inherited from Middle Chinese, \emph{syuw}, \emph{ngin} and \emph{doj} (Guangyun, in Baxter's transcription). The dictionary glosses the last as a raised square platform of earth.
\end{cousinweb}

\begin{grammarlens}[title={What you've built}]
One more word for asking for what I want.
\end{grammarlens}

\subsection*{Your turn}
\begin{itemize}
\raggedright
  \item \emph{Say it:} \emph{shōuyíntái}
  \item \emph{Say it:} \emph{shōuyíntái}, once more
  \item \emph{Say it:} \emph{shōuyíntái zài ménkǒu} — The checkout is by the door.
  \item \emph{Recall:} say the Mandarin for to take goods back, then the Mandarin for customer, then say \emph{shōuyíntái} again
\end{itemize}

\begin{tcolorbox}[breakable,colback=teal!4,colframe=teal!35!black,title={Before you move on}]
What does shōuyíntái mean? (\textquotedblleft{}Checkout; till\textquotedblright{}.) Say it once more.
\end{tcolorbox}
